
\section{Limitations and Conclusion}

\subsection{Limitations}

This study has several limitations. First, the pilot
evaluation uses only 10 episodes, with six unsafe
episodes in the test set. This sample is too small
to support broad conclusions about monitoring
effectiveness.

Second, the experiment uses a keyword-based monitoring
baseline. Its detection performance may depend on
the selected keywords and the wording of the examples.
Therefore, the results should not be generalized to
more advanced LLM-based monitors.

Third, only four categories of unsafe behavior are
considered: reward hacking, deception, sabotage, and
harmful compliance. Other forms of unsafe behavior
are outside the scope of the current dataset.

Finally, the load-testing results were obtained under
a specific benchmark configuration. They do not
necessarily represent performance in real-world
production environments.

\subsection{Conclusion}

This study investigated whether access to reasoning
traces can improve the detection of unsafe behavior
in large language models. Three research questions
guided the investigation.

For RQ1, the category-level statistical analysis was
degenerate and did not provide sufficient evidence
to establish differences between categories.

For RQ2, the observed recall in the CoT-plus-output
condition was 1.00. However, the binomial test did
not establish that recall was significantly greater
than the target of 0.70 at the 0.05 significance
level.

For RQ3, the exact McNemar test indicated a
statistically significant difference between the
paired monitoring conditions (p = 0.015625).
In this pilot experiment, access to reasoning traces
was associated with improved detection compared
with output-only monitoring.

Overall, the findings provide preliminary evidence
that reasoning traces may contain useful signals
for monitoring unsafe behavior. However, larger
datasets and more diverse monitoring systems are
needed to validate this observation.

\subsection{Future Work}

Future work will focus on three directions. First,
the evaluation dataset should be expanded to include
more safe and unsafe examples across all four
categories. Second, the keyword-based baseline
should be compared with LLM-based monitoring
approaches. Third, additional experiments should
evaluate robustness, false-positive rates, and
latency under more realistic workloads.

These improvements will help determine whether
reasoning-trace access provides consistent benefits
for AI safety monitoring across different models
and evaluation settings.
